\section{紧凑等价检查与 MLIR 展望}
\label{sec:validation}

\subsection{行为检查：支持到哪里，结论就说到哪里}

验证把“标准输出中的一个整数加换行，退出状态为 0”作为可观察契约。六个输入覆盖负数下界、零次循环、最短正输入、多次迭代、上界和超上界。表~\ref{tab:tests} 中负数先夹到 0，所以最终输出 $0!=1$；“夹到 0”不等于“直接输出 0”。C 形态与 SysY 采用同一算法，表中只把后者列作源语义代表。

\begin{table}[tb]
\caption{六个输入的期望值与三种作者表示的结果。}
\label{tab:tests}
\centering\small
\begin{tabular}{@{}rrlll@{}}
\toprule
输入 & 期望输出 & SysY 源 & 作者 IR & 作者汇编\\
\midrule
$-3$ & 1 & 通过 & 通过 & 通过\\
0 & 1 & 通过 & 通过 & 通过\\
1 & 1 & 通过 & 通过 & 通过\\
5 & 120 & 通过 & 通过 & 通过\\
10 & 3628800 & 通过 & 通过 & 通过\\
20 & 3628800 & 通过 & 通过 & 通过\\
\midrule
合计 & -- & 6/6 & 6/6 & 6/6\\
\bottomrule
\end{tabular}
\end{table}

这项检查能发现常见的边界、循环、调用和输出差错，并支持“在声明的六个输入上可观察行为一致”。它不是形式等价证明：输入空间没有穷尽，工具链和环境也是前提。更强结论需要全域符号推理或经证明的翻译验证，并明确整数、I/O 与终止语义。版本漂移也会改变打印 IR、优化流水线和指令选择；因此本文把版本、目标三元组、ISA/ABI 和计数口径列入附录，并把所有大小称为本次实验观测。

\subsection{展望：MLIR 把“何时丢弃抽象”变成设计对象}

LLVM IR 适合显式控制流、内存与机器式标量操作，但领域操作若过早拆散，后续领域变换可能无法恢复其结构。MLIR 提出可扩展方言（dialect）与多级中间表示，使领域、算法、内存和硬件抽象在同一基础设施中逐步降低\cite{lattner-et-al-mlir-2021}。本阶乘并不需要 MLIR 才能高效实现；这里只用它解释机制，不声称本地运行或加速器验证。

一个概念路径是：\code{func+arith+scf} 保留函数、整数运算和结构化循环；降低到 \code{func+arith+cf} 后，循环成为块、分支与块参数；再转换到 LLVM 方言并翻译为 LLVM IR，复用既有优化和 RV64 后端。方言转换框架用转换目标声明哪些操作合法，用重写模式替换非法操作，并可用类型转换器桥接暂时混合的类型；部分转换允许不同抽象共存\cite{mlir-dialect-conversion-2026,mlir-llvm-target-2026,mlir-toy-lowering-2026}。

\begin{center}
\small
\fbox{\code{scf} 结构化循环 $\rightarrow$ \code{cf} 块/分支 $\rightarrow$ LLVM 方言 $\rightarrow$ LLVM IR $\rightarrow$ RV64}
\end{center}

课程进阶材料点名的 AscendNPU IR VecAdd 展示了同一思想在加速器上的形态：官方 HIVM 示例把设备入口、全局内存（\code{gm}）到统一缓冲区（\code{ub}）的搬运、\code{hivm.hir.vadd} 以及写回保持为显式操作；架构说明再把 HFusion、HIVM、低层 MLIR 与 LLVM IR/算子二进制组织成逐级降低关系\cite{ascend-vecadd-2026,ascend-architecture-2026}。本文没有执行 CANN/BiSheng 专用编译器或 NPU 真机，因此这只是文档支撑的设计展望，不是本阶乘实例已经经过的阶段，也不是设备性能实验。

其价值是\emph{表示时机}：在最后一个需要高层结构的变换完成后再降低。代价同样真实：每条方言边界都需要合法性规则、重写、类型桥接、诊断和版本维护；方言也不会自动组合，保留结构更不保证程序更快。因此，MLIR 更适合确有多个领域或硬件层级、而 LLVM IR 会过早抹平关键信息的系统，不应仅因“层次更多”就用于小型标量编译器。
